\section{A bilinear frame estimate without a scale loss}
\label{sec:frame}

This section supplies the estimate for long, nearly flat frequency blocks.
Its inputs need not solve the Schr\"odinger equation: their profiles may
vary smoothly on the packet time scale.  The proof successively merges
frequency caps and removes large pieces.  A positive energy recurrence
records the removals. Low-frequency flux bounds and a quartic potential
control its errors uniformly in the number of merges.

The cap-merging and amplitude-pruning scheme is related to the high--low
method of Guth--Maldague--Wang~\cite[Sections~2.2--2.4 and~3.3]{GMW24}
and its $p$-adic implementation by Guo--Li--Yung~\cite[Section~7]{GLY23}.
Retaining discarded quadratic energy and controlling the resulting
square function by quartic telescoping are central in
Deng--Fan--Guo--Guo--Luo~\cite[Sections~4.2 and~6.2]{DFGGL26}
and the related contemporary work of
Li--Yung~\cite[Sections~5.2 and~5.5]{LiYung26}.
Those latter arguments exploit exact truncation and Fourier-support
identities. Here smooth real-variable masks introduce frequency leakage;
the low-frequency flux estimate and Gaussian quartic potential below
control the resulting errors uniformly in the number of scales.
The smooth-mask, energy-flux and quartic-potential construction and the
scale-uniform estimate are proved locally; they are not conclusions imported
from the Guth--Maldague--Wang decoupling theorem.

\begin{theorem}[Bilinear frame estimate]\label{thm:frame}
Fix bounded frequency and profile boxes, a bound on phase-space
multiplicity, and a sufficiently large envelope exponent $P$.
There is a fixed finite number of profile derivatives for which the
following holds uniformly.  Let $h\geq1$ and
$h^{-1/2}\leq\delta\lesssim1$.  For $i=1,2$, let
\[
 V_i(x,t)=\sum_{\ell\in\mathcal I_i}s_\ell
 e^{i(x\xi_\ell-t\xi_\ell^2)}
 H_\ell\left(\frac{x-b_\ell-2\xi_\ell t}{h},\frac{t}{h^2}\right)
\]
be a finite packet sum.  The profiles are supported in the fixed box
and have the stipulated uniform derivative bounds.  The centers $b_\ell$
may be arbitrary real numbers, and the pairs
$(b_\ell/h,h\xi_\ell)$ have bounded multiplicity in every unit
phase-space ball.  The frequencies of the two fields lie in intervals
of length $\delta$ with distance at least a fixed positive multiple of
$\delta$.  Put
\[
 S_i(x,t)^2=\sum_{\ell\in\mathcal I_i}|s_\ell|^2
 \left(1+\frac{|x-b_\ell-2\xi_\ell t|}{h}
             +\frac{|t|}{h^2}\right)^{-P},
 \qquad H_i^*=h^3\sum_{\ell\in\mathcal I_i}|s_\ell|^2.
\]
For $A_i,B_i>0$, write $\alpha_i=A_i/B_i$.  If
$\alpha_i\geq C$, for a sufficiently large fixed $C$, the number of
unit lattice squares containing a point where simultaneously
$|V_i|\geq A_i$ and $S_i\leq B_i$ for $i=1,2$ satisfies
\begin{equation}\label{eq:src8}
 \#\leq C'\delta^{-C'}\alpha_1^{-2}\alpha_2^{-2}
 \left(\frac{H_1^*}{A_1^2}+\frac{H_2^*}{A_2^2}\right).
\end{equation}
The constants depend only on the fixed bounds just specified.  In
particular, they are independent of $h$, the number of packets, and the
number of scales in the proof.  The loss in the separation parameter
has at most a fixed polynomial power.
\end{theorem}

We always take one witness per counted unit square; the witnesses for
all four inequalities in a square must be the same point.  All norms
below, unless a variable is specified, are on $(x,t)\in\R^2$.  Spatial
Fourier multipliers act only in $x$.  Fourier normalization constants
in positive potentials are absorbed into their definitions.  Set
\[
 \mathcal L=i\partial_t+\partial_x^2,
 \qquad D_a=\partial_t+2a\partial_x-ia^2.
\]
The latter derivative annihilates the carrier at frequency $a$; useful
identities are
\begin{equation}\label{eq:frame-derivatives}
 \mathcal L=iD_a+(\partial_x-ia)^2,
 \qquad
 D_a-D_b=2(a-b)(\partial_x-ib)-i(a-b)^2.
\end{equation}

\subsection{Reduction to small, exactly localized frequency pieces}

\begin{lemma}[Frequency layers and initial normalization]
\label{lem:frame-initial}
It suffices to prove \eqref{eq:src8} for packets whose spatial
frequencies are supported in caps of width $O(h^{-1})$ about their
assigned centers, whose spatial profiles have sufficiently rapid
polynomial decay, and whose native time supports remain bounded.
After dividing a field by its envelope bound $B$, its witness level is
$\alpha=A/B$ and its envelope is at most one.  Given
\[
 \lambda=C_*/\alpha\leq1,
\]
deleting all packets of normalized amplitude greater than $\lambda$
changes the field at every witness by at most $C\lambda^{-1}$.
The remaining field $g_0$ has an energy parameter
$D\lesssim H^*/B^2$ such that $\|g_0\|_2^2\leq D$.  At a dyadic
scale $\Delta_0\asymp h^{-1}$, its fine cap pieces $g_{0,b}$ satisfy,
for all fixed orders needed below,
\begin{align}
 \| (\partial_x-ib)^sD_b^r g_{0,b}\|_\infty
 &\leq C_{r,s}\lambda\Delta_0^{s+2r},
 \label{eq:frame-initial-sup}\\
 \sum_b\|(\partial_x-ib)^sD_b^r g_{0,b}\|_2^2
 &\leq C_{r,s}\Delta_0^{2s+4r}D.
 \label{eq:frame-initial-square}
\end{align}
In particular $\|\mathcal L^r g_0\|_2\lesssim_r
\Delta_0^{2r}\sqrt D$.
\end{lemma}

\begin{proof}
Smoothly partition the spatial Fourier transform of each native profile
into layers $k+[-2,2]$, $k\in\mathbb Z$, writing
$\langle k\rangle=1+|k|$.  A fixed, sufficiently large number of
derivatives gives any prescribed fixed polynomial layer decay together
with the finitely many spatial decay and derivative estimates that we
will use.  Assign the new carrier center $\xi_\ell+k/h$.
In its native coordinates the profile undergoes a translation by
$2k$ times native time and acquires a temporal oscillation of frequency
$k^2$.  Every required derivative and envelope comparison therefore
costs only a fixed polynomial in $\langle k\rangle$.  Choose the
original decay orders larger than these costs and put the remaining
decay into the layer amplitudes.  A common frequency shift preserves
the phase-space counts.

At a witness the sum of all layers with $|k|\geq c\delta h$ is bounded
by
\[
 C_N(\delta h)^{-N}\sqrt h\,B.
\]
Indeed, in each original frequency bin of width $h^{-1}$, the spatial
tails sum uniformly by bounded offset counts; weighted
Cauchy--Schwarz over the $O(h)$ frequency bins gives $\sqrt h$.
Since $\delta h\geq\sqrt h$, the displayed error is negligible for
large $h$.  Bounded scales are vacuous once the required ratio constant
is large, by the same envelope Cauchy--Schwarz estimate.
Choose the fixed $c>0$ small enough that the retained frequency layers
still lie in separated intervals of length $O(\delta)$.

If the retained sum is large, some layer has modulus at least
$c'\langle k\rangle^{-2}A$.  For any prescribed fixed large $N$ its
envelope bound and energy can be taken to be
$C_N\langle k\rangle^{-N}B$ and
$C_N\langle k\rangle^{-2N}H^*$, respectively.  Thus the ratio condition
persists, and applying the desired estimate at these thresholds and
summing the two layer indices converges.  For example, the ratio
factors gain $\langle k\rangle^{-2(N-2)}$ and the corresponding
energy-to-level ratio gains the same power.  Constants and an initial
finite set of layers are absorbed into the fixed ratio threshold.

We may consequently work with one exactly cap-supported layer, divided
by $B$.  If $|s_\ell|>\lambda$, then
$|s_\ell|\leq\lambda^{-1}|s_\ell|^2$.  The pointwise profile decay,
chosen stronger than the envelope decay, bounds the deleted sum by
$C\lambda^{-1}S^2\leq C\lambda^{-1}$ at witnesses.

For a fixed fine frequency bin, transported spatial intercepts still
have bounded counts per interval of length $h$ on $|t|\lesssim h^2$:
velocity variation inside that bin changes positions by only $O(h)$.
Summing profile tails over these offsets gives
\eqref{eq:frame-initial-sup}, since all retained amplitudes are at most
$\lambda$.  Each derivative $D_b$ costs $\Delta_0^2$, and each
$\partial_x-ib$ costs $\Delta_0$, including the harmless change from
an individual center to the bin center.  Weighted Cauchy--Schwarz
within a bin gives \eqref{eq:frame-initial-square}; a single packet
has squared spacetime measure scale $h^3$.  Spatial frequencies of
distinct bins have bounded overlap.  Enlarging $D$ by a fixed factor
proves all the stated assertions.
\end{proof}

\subsection{A positive frame iteration and smooth deletion masks}

Choose a large fixed dyadic integer $K$, whose size will be fixed in
Lemma~\ref{lem:frame-masks}.  Set $\Delta_j=K^j\Delta_0$, ending at
$\Delta_n\asymp_K c_0\delta$, where $c_0>0$ is a sufficiently small
fixed constant.  A shorter range with only boundedly many fine caps
per frequency interval again makes large ratios impossible.  The
terminal cap size is chosen small enough to preserve the frequency
gap under all the enlargements below.

Fix a real $p\in C_c^\infty(\R)$ with $\int p^2=1$.  At stage $j$,
write $\Delta=\Delta_j$ and
\[
 P_a=p((\xi-a)/\Delta),\qquad
 \int_a=\int_\R\frac{da}{\Delta},\qquad
 \int_aP_a^2=I.
\]
With $g=g_{j-1}$ define
\begin{equation}\label{eq:frame-step}
 \begin{gathered}
 c_a=P_ag,\qquad b_a=\rho_ac_a,\qquad d_a=c_a-b_a,\\
 g_j=\int_aP_ab_a,\qquad R_j=g_{j-1}-g_j=\int_aP_ad_a,\\
 L_a=|c_a|^2-|b_a|^2,\qquad
 L_j=\int_aL_a,\qquad l_j=\int_{\R^2}L_j.
 \end{gathered}
\end{equation}
The masks will obey $0\leq\rho_a\leq1$.  Analysis by the $P_a$ is an
isometry and synthesis is its adjoint, hence a contraction.  Therefore
\begin{equation}\label{eq:frame-contraction}
 \|g_j\|_2^2\leq\|g_{j-1}\|_2^2-l_j,\qquad
 \sum_jl_j\leq D,\qquad
 \int_a\|d_a\|_2^2\leq l_j.
\end{equation}
Here $|d_a|^2\leq L_a$ follows from
$(1-\rho_a)^2\leq1-\rho_a^2$.
The outer projection in synthesis ensures that each step inflates
spatial frequency support by only $O(\Delta_j)$, even though masking
itself creates frequencies.  The total inflation is
$O(\sum_j\Delta_j)=O(\Delta_n)$, preserving the gap.
All states retain the original bounded time support.

We also fix positive spatial convolutions $W_j$ with kernels comparable
to
$\Delta_j(1+\Delta_j|x|)^{-P_0}$, where $P_0>P$ is a sufficiently
large fixed even integer.  Their common rescaled multiplier $w$ is
even, compactly supported, sufficiently smooth, and satisfies $w(0)=1$.
One construction takes the normalized autocorrelation of a high
convolution power of $e^{-x}\mathbf1_{[0,1]}(x)$ as multiplier.
The Fourier modulus of the truncated exponential is comparable to
$(1+|\xi|)^{-1}$, so the resulting positive kernel has the asserted
two-sided polynomial bounds.  Increasing that fixed convolution power
supplies the required multiplier smoothness.

At every normalized witness $z_0$,
\begin{equation}\label{eq:src14}
 W_0|g_0|^2(z)\lesssim(1+|z-z_0|/h)^P,
\end{equation}
with a harmless enlargement of the fixed envelope powers if needed.
To verify this, compact support of the multiplier of $W_0$ eliminates
products of nonneighboring fine caps.  Bound each remaining product
by the sum of the squared cap moduli, and then use weighted
Cauchy--Schwarz and the transported offset counts within each cap.
Positive convolution preserves the resulting sum of squared packet
amplitudes with decaying spatial tails.  Bounded velocities give its
polynomial comparison with the envelope at $z_0$; the original native
time support is bounded.  Choosing all decay powers in advance larger
than $P$ proves \eqref{eq:src14}.

The global inequality $\sum_jl_j\leq D$ does not by itself control
how much field is removed at a particular witness.  The masks below will
satisfy the pointwise bound $|R_j|\lesssim_K\lambda^{-1}W_jL_j$.
We therefore need a positive recurrence for the smoothed losses $W_jL_j$.
The following identity holds for any masks $0\leq\rho_a\leq1$ in
\eqref{eq:frame-step}; it does not use their construction.

Define
\[
 S_j^0=W_j|g_j|^2,\qquad
 B_j^0=\int_a|d_a-P_aR_j|^2\geq0.
\]
The term $B_j^0$ compares the removed cap family $d_a$ with the
analysis $P_aR_j$ of its synthesized field.  Define the two error fields
and their accumulated sum by
\[
 H_j=(W_j-W_{j-1})|g_{j-1}|^2,\qquad
 \Phi_j=|g_j|^2-|g_{j-1}|^2+L_j+B_j^0,
 \qquad Y_k=\sum_{j\leq k}(H_j+W_j\Phi_j).
\]
The exact telescoping identity is
\begin{equation}\label{eq:frame-recurrence}
 S_k^0+\sum_{j\leq k}W_j(L_j+B_j^0)=S_0^0+Y_k.
\end{equation}
The term $H_j$ records the change of smoothing scale; $\Phi_j$ records
the local error in passing between cap analysis and synthesis.
Every term on the left is nonnegative.  In particular, the witness
baseline \eqref{eq:src14} gives
$\sum_{j\leq k}W_jL_j\lesssim1+|Y_k|$ at a witness.
Thus the accumulated deletions remain controlled until $Y_k$ becomes
large.  We will prove the uniform maximal estimate
\begin{equation}\label{eq:src19}
 \left\|\sup_{0\leq k\leq n}|Y_k|\right\|_2
 \leq C_K\lambda\sqrt D,
 \qquad Y_0=0.
\end{equation}

\begin{lemma}[Masks, persistence, and loss bounds]\label{lem:frame-masks}
The masks can be chosen so that
\begin{equation}\label{eq:src15}
 |b_a|\leq\lambda,\qquad
 \|\mathcal L P_ab_a\|_\infty\leq C_1\lambda\Delta^2,
\end{equation}
where $C_1$ is independent of the large fixed $K$.  For every required
fixed derivative order,
\begin{equation}\label{eq:src16}
 \|D_a^rP_ab_a\|_\infty\leq C_{r,K}\lambda\Delta^{2r},
 \qquad
 \int_a\|D_a^rP_ab_a\|_2^2\leq C_{r,K}\Delta^{4r}D.
\end{equation}
The removed pieces obey
\begin{equation}\label{eq:src17}
 \begin{split}
 \|(\partial_x-ia)^sD_a^rd_a\|_1
 &\leq C_{K,r,s}\Delta^{s+2r}\lambda^{-1}\int L_a,\\
 \|(\partial_x-ia)^sD_a^rd_a\|_2^2
 &\leq C_{K,r,s}\Delta^{2s+4r}\int L_a,
 \end{split}
\end{equation}
and, pointwise,
\begin{equation}\label{eq:src18}
 |R_j|\leq C_K\lambda^{-1}W_jL_j.
\end{equation}
All derivative constants are uniform in the stage number.
\end{lemma}

\begin{proof}
\emph{Construction of the masks.}
Summing the preceding caps that meet $P_a$ gives
\begin{equation}\label{eq:frame-incoming}
 \|c_a\|_\infty\leq CK\lambda,\qquad
 \|\mathcal Lc_a\|_\infty
 \leq CC_1\lambda\Delta^2/K.
\end{equation}
There are $O(K)$ preceding normalized centers, and their curvature
scale is $(\Delta/K)^2$.  The initial bounds give the same estimates
at the first merge.

Take a nonnegative smooth partition of unity
$\sum_e\chi_e(\Delta^2t)=1$ by translates with bounded support and
overlap, each bounded below on a fixed inner interval.  In the
co-moving coordinate $r_x=\Delta(x-2at)$, mark every unit cell on which
$|c_a|>\lambda$ somewhere during the open positive support of the
$e$th time cutoff, and let $A_e$ be their union.  With a fixed
nonnegative $\psi\in C_c^\infty((-1,1))$ of integral one, take
\[
 m_e=\psi*\mathbf1_{\{r:\dist(r,A_e)<K+2\}}.
\]
For $K\geq4$, this function is one within distance $K$ of $A_e$,
zero outside its $2K$ enlargement, and lies in $[0,1]$; its fixed-order
derivative bounds are independent of $K$.  Define
\[
 1-\rho_a=\sum_e\chi_e(\Delta^2t)m_e(r_x).
\]
Choose the time cutoffs flat to all required orders at the endpoints
of their positive supports.  Marking can be tested on dense countable
sets in the open slabs, using continuity, so the construction is
measurable also in $a$.  Half-open cell conventions, or fixed slight
enlargements, leave all bounds unchanged.  At a point with
$|c_a|>\lambda$ every active index is marked, giving $\rho_a=0$ and
therefore $|b_a|\leq\lambda$.

\emph{The curvature bound.}
Where a derivative of $\rho_a$ is nonzero, there is no same-time peak
within co-moving distance $K/2$.  Otherwise every active index would
be identically fully marked locally, and all derivatives of the mask
would vanish.  Local Bernstein, using $|c_a|\leq\lambda$ in this
neighborhood and the global bound $CK\lambda$ against rapidly
decreasing reproducing-kernel tails outside it, gives
$|(\partial_x-ia)c_a|\lesssim\Delta\lambda$ there.  Now
\[
 \mathcal L(\rho_ac_a)=\rho_a\mathcal Lc_a+
 \bigl(i(\partial_t+2a\partial_x)\rho_a+
                  \partial_x^2\rho_a\bigr)c_a
 +2(\partial_x\rho_a)(\partial_x-ia)c_a.
\]
Scaled derivatives of the masks cost $\Delta$ spatially and
$\Delta^2$ in co-moving time.  Thus
$\|\mathcal Lb_a\|_\infty\lesssim
\lambda\Delta^2(1+C_1/K)$.
The outer projection has bounded $L^1$ kernel.  Choose $C_1$ first
sufficiently large independently of $K$, and then $K$ sufficiently
large, to close \eqref{eq:src15}.

\emph{Uniform higher derivatives.}
Order zero follows from bounded masks and frame contraction; spatial
derivatives of projected pieces follow from Bernstein.  For higher
$D_a$ orders use \eqref{eq:frame-derivatives} to express the incoming
piece in preceding centers $\widetilde a$.  After normalization by
$\Delta^{2r}$, the coefficient of the preceding top-order supremum
bound is at most $CK^{1-2r}$.  In integrated center $L^2$ norm it is
at most $CK^{1/2-2r}$: Cauchy--Schwarz in the preceding-center window
costs $O(\sqrt K)$, while the current-center window for each fixed
preceding center has $da/\Delta$ measure $O(1)$.
Every other term in the center change uses a lower affine time order;
its spatial derivatives cost the preceding Bernstein scale.
In $D_a^r(\rho_ac_a)$ the top term is $\rho_aD_a^rc_a$, with
$|\rho_a|\leq1$, and all remaining terms use lower orders and
controlled co-moving derivatives of $\rho_a$.
Fix the maximum required derivative order in advance.  The top
coefficients are contractive for every $r\geq1$ in this finite list
after increasing $K$ once.  This enlargement leaves $C_1$ unchanged
and only improves both the curvature closure and the Duhamel source
smallness.  Induction on $r$, and then on stages with that
contraction, proves \eqref{eq:src16}.  The first step uses the fine
lattice-bin estimates.  The same reasoning gives mixed spatial and
affine supremum bounds for $c_a,d_a$, with constants depending on $K$.

\emph{Persistence of a marked peak.}
A peak in a marked cell forces spatial squared mass at least
$c\lambda^2/\Delta$ within co-moving distance $K/2$ of that cell at
\emph{every} comparison time in its time slab.  To see this, propagate
from the comparison time $t$ to the peak time $t_0$ by Duhamel, inserting
a smooth cap cutoff equal to one on the spatial Fourier supports of both
$c_a$ and $\mathcal Lc_a$.  Since $|t-t_0|=O(\Delta^{-2})$, the
co-moving propagator kernels are uniformly rapidly decreasing at
scale $\Delta^{-1}$, with $L^1$ norm $O(1)$ and $L^2$ norm
$O(\sqrt\Delta)$.  The source contribution is at most
$CC_1\lambda/K$ by \eqref{eq:frame-incoming}.  The main term outside
distance $K/2$ is negligible by its kernel tails and the global
$CK\lambda$ bound.  With $K$ fixed sufficiently large, the remaining
main term has modulus at least $c\lambda$.  Cauchy--Schwarz with its
$O(\sqrt\Delta)$ kernel norm proves the asserted mass lower bound.
This argument explicitly permits nonzero $\mathcal Lc_a$.

\emph{Charging the activated sets.}
We have $L_a\geq\sum_e\chi_em_e|c_a|^2$.  For each marked cell and
slab, its enlarged activated region has spacetime volume
$O_K(\Delta^{-3})$.  Integrating the preceding persistence bound over
the inner times where $\chi_e$ is bounded below gives a charge
$c\lambda^2\Delta^{-3}$ in that cell's fully marked neighborhood.
These neighborhoods have overlap $O_K(1)$, including across slabs.
Consequently the total activated volume, including the supports of
mask derivatives counted slab by slab, is at most
$C_K\lambda^{-2}\int L_a$.
The mixed derivative supremum bounds for $d_a$ are
$C_{K,r,s}\lambda\Delta^{s+2r}$ on this activated set and vanish
off it.  Multiplying by its volume proves both inequalities in
\eqref{eq:src17}.

Finally fix a time with $\chi_e>0$.  The spatial projection of
$\chi_em_ec_a$ is bounded by summing its rapidly decreasing kernel
over activated cells, using $|c_a|\leq CK\lambda$.
Persistence gives at that same time a lower squared-mass bound in
each corresponding fully marked neighborhood.  Since the positive
kernel of $W_j$ decays more slowly than the projection kernel, and
the enlarged neighborhoods have bounded overlap, this gives
\[
 |P_a(\chi_em_ec_a)|
 \leq C_K\lambda^{-1}W_j(\chi_em_e|c_a|^2).
\]
The common factor $\chi_e(t)$ causes no difficulty even near slab
endpoints.  More explicitly, with $u=\Delta(x-2at)$ and marked
cell indices $q$, the kernel upper bound is
$C_K\lambda\sum_q(1+|u-q|)^{-P_0}$.
Persistence, kernel comparability on the fully marked neighborhoods,
and their $O_K(1)$ overlap give
\[
 \lambda^2\sum_q(1+|u-q|)^{-P_0}
 \lesssim_K W_j(m_e|c_a|^2)(x,t).
\]
Multiplying by $\chi_e(t)$, summing $e$, and integrating $a$ proves
\eqref{eq:src18}.
\end{proof}

\subsection{Energy flux and the quartic potential}

We now bound the two errors in \eqref{eq:frame-recurrence}.
The deletion flux $W_j\Phi_j$ will be charged to $l_j$; the changes
$H_j$ of smoothing scale require a quartic potential.

\begin{lemma}[Low-frequency flux]\label{lem:frame-flux}
For a smooth spatial projection $\operatorname{Proj}_r$ to
$|\xi|\asymp r$, where $r\lesssim\Delta_j$,
\begin{equation}\label{eq:src20}
 \|\operatorname{Proj}_rW_j\Phi_j\|_2^2
 \leq C_K\lambda^2(r/\Delta_j)l_j.
\end{equation}
\end{lemma}

\begin{proof}
Write $g=g_{j-1}$, $R=R_j$, and $\Delta=\Delta_j$.
Expanding $B_j^0$ and using $g_j=g-R$ gives
\begin{align*}
 \Phi_j={}&-2\operatorname{Re}\left(g\overline R-
                      \int_a(P_ag)\overline{d_a}\right)\\
 &+2\operatorname{Re}\left(|R|^2-
                      \int_ad_a\overline{P_aR}\right)
       +\int_a|P_aR|^2-|R|^2.
\end{align*}
For either of the first two differences the product multiplier
contains, up to sign,
$p((\xi_1-a)/\Delta)-p((\xi_2-a)/\Delta)$.
On the output band $|\xi_1-\xi_2|\asymp r$ it costs $O(r/\Delta)$.
Both inputs may be restricted to $O(\Delta)$ about $a$: the cutoff
on one factor and the output support force the same localization on
the other, including the otherwise unprojected $d_a$.
The last difference has symbol
\[
 \int_\R p((\xi_1-a)/\Delta)p((\xi_2-a)/\Delta)
                 \frac{da}{\Delta}-1,
\]
which vanishes at $\xi_1=\xi_2$ and has the same bound.
After bin and band localization, all fixed rescaled derivatives in
$(\xi_1-a)/\Delta$ and $(\xi_1-\xi_2)/r$ satisfy that bound.
Fourier expansion therefore represents each multiplier by spatially
translated products with total coefficient or kernel mass
$O_K(r/\Delta)$.

Here is the loss-sensitive norm bound for those products.
Let $Q_b$ be smooth lattice projections of width $\Delta$, allowing
fixed enlargements.  From \eqref{eq:src16} and the center-change
identity,
\[
 \|D_b^sQ_bg\|_\infty+
 \|D_b^sQ_bR\|_\infty\leq C_{s,K}\lambda\Delta^{2s}.
\]
Moreover,
\begin{equation}\label{eq:frame-loss-bins}
 \sum_b\|D_b^sQ_bR\|_2^2
 \leq C_{s,K}\Delta^{4s}l_j.
\end{equation}
Indeed $Q_bR=\int_aQ_bP_ad_a$ uses only $|a-b|\lesssim\Delta$;
apply Cauchy--Schwarz in this normalized window and
\eqref{eq:src17}, observing boundedly many $b$ per $a$.
The terms with an explicit $d_a$ are handled in the same way after
grouping the continuous centers into lattice bins and inserting the
permitted spatial projections.  All these estimates persist under
the spatial translations in the multiplier expansion.  Indeed, for
$T_yf(x,t)=f(x-y,t)$, one has $D_bT_y=T_yD_b$; translations preserve
both spatial cap support and the global $L^\infty,L^2$ norms used
here.  Thus even translations of size $r^{-1}$ in the output-band
variable have no cost depending on $\Delta/r$.

Thus the output from a bin centered at $b$ has $L^2$ norm at most
$C_K\lambda(r/\Delta)l_b^{1/2}$, with
$\sum_bl_b\lesssim_K l_j$.  It has the same bound after any of the
fixed normalized product derivatives
$((\partial_t+2b\partial_x)/\Delta^2)^s$ that we need.
On a product with a conjugated factor this derivative splits as
$D_b$ on the factors, with the carrier constants canceling.
On spacetime Fourier variables $(\xi,\tau)$, the weighted overlap
of the bin strips obeys
\[
 \sum_b\left(1+\frac{|\tau+2b\xi|}{\Delta^2}\right)^{-4}
 \lesssim1+\frac\Delta r.
\]
The centers $b$ have spacing $\Delta$, so their strip centers in
$\tau$ have spacing comparable to $\Delta r$ and width
$\Delta^2$.  Weighted Cauchy--Schwarz, followed by Plancherel and
the differentiated product bounds, gives
\[
 \|\operatorname{Proj}_rW_j\Phi_j\|_2^2
 \lesssim_K\lambda^2(r/\Delta)^2(1+\Delta/r)\sum_bl_b,
\]
which is \eqref{eq:src20}.
\end{proof}

Flux alone does not control the change in the spatial smoothing
scale.  We obtain that control from a positive quartic expression,
whose Gaussian multiplier permits an explicit calculation.

\begin{lemma}[Quartic control of the growth bands]
\label{lem:frame-quartic}
Let $G_j(u)=\exp(-cu^2/\Delta_j^2)$, for a fixed $c>0$.
Then
\begin{equation}\label{eq:src22}
 \sum_{j=1}^n\int_{\R^2}(G_j-G_{j-1})(u)
       |\widehat{|g_{j-1}|^2}(u,\tau)|^2\,du\,d\tau
 \lesssim_K\lambda^2D.
\end{equation}
\end{lemma}

\begin{proof}
For a spatial Gaussian multiplier $q_{j,a}$ of width $\Delta_j$
centered at $a$, define the normalized potential
\[
 \mathcal P_j(f)=\int_\R\int_{\R^2}|q_{j,a}f(z)|^4\,dz
                                      \frac{da}{\Delta_j}.
\]
On quartets with alternating sums of both spatial and temporal
frequencies equal to zero, put
$u=\xi_1-\xi_2$ and $v=\xi_1-\xi_4$.
The common midpoint is
$(\xi_1+\xi_3)/2=(\xi_2+\xi_4)/2$; the sum of the four squared
deviations from it is $u^2+v^2$.  Integrating the Gaussian center
therefore gives exactly the multiplier $G_j(u)G_j(v)$, after fixed
normalization.  Expanding this product yields
\begin{equation}\label{eq:frame-potential-increment}
 \mathcal P_j(g)-\mathcal P_{j-1}(g)
 =2\int(G_j-G_{j-1})(u)
                |\widehat{|g|^2}(u,\tau)|^2\,du\,d\tau
    +\operatorname{Err}_j,
\end{equation}
where $g=g_{j-1}$ and the error multiplier is
\[
 E_j(u,v)=(1-G_j(u))(1-G_j(v))
                 -(1-G_{j-1}(u))(1-G_{j-1}(v)).
\]
The two quadratic terms in the expansion are equal by interchanging
$u$ and $v$.

We establish the quantitative error estimate
\begin{equation}\label{eq:src21}
 |\operatorname{Err}_j|\lesssim_K\lambda^2
 \left[\left(\frac{\Delta_0}{\Delta_{j-1}}\right)^2D+
 \sum_{m<j}\left(\frac{\Delta_m}{\Delta_{j-1}}\right)^2l_m\right].
\end{equation}
Write $\Delta'=\Delta_{j-1}$ and partition every input into smooth
spatial caps $Q_b$ of width $\Delta'$.  The quartet identity
\begin{equation}\label{eq:frame-quartet}
 2uv=\xi_1^2-\xi_2^2+\xi_3^2-\xi_4^2
     =\sum_{\nu=1}^4(-1)^{\nu+1}(\tau_\nu+\xi_\nu^2)
\end{equation}
allows division of $E_j$ by $2uv/(\Delta')^2$, placing a normalized
curvature derivative $\mathcal L/(\Delta')^2$ on one input.
The quotient is smooth at $u=0$ or $v=0$, since $E_j$ vanishes
quadratically in each variable.  The sign of the multiplier
$\mathcal L$, namely $-(\tau+\xi^2)$, is irrelevant to the bounds.

\emph{Norm bookkeeping for a distinguished derivative.}
On the input that receives this first derivative write
$g=g_0-\sum_{m<j}R_m$.  For every required fixed $q\geq1$,
\begin{align}
 \left(\sum_b\|Q_b\mathcal L^qg_0/(\Delta')^{2q}\|_2^2\right)^{1/2}
 &\lesssim_K(\Delta_0/\Delta')^{2q}\sqrt D,
 \label{eq:frame-curvature-base}\\
 \sum_b\|Q_b\mathcal L^qR_m/(\Delta')^{2q}\|_1
 &\lesssim_K(\Delta_m/\Delta')^{2q}\lambda^{-1}l_m.
 \label{eq:frame-curvature-loss}
\end{align}
The first assertion is the initial mixed-derivative estimate and
bounded spatial Fourier overlap.  For the second use
$R_m=\int_aP_ad_a$ at scale $\Delta_m$ and
$\mathcal L=iD_a+(\partial_x-ia)^2$.
Each derivative costs $\Delta_m^2$ in the $L^1$ bound
\eqref{eq:src17}.  The projection kernels have bounded $L^1$ norm,
and each earlier synthesis center meets only boundedly many current
bins, proving \eqref{eq:frame-curvature-loss}.
On every other input keep $g=g_{j-1}$ itself.  Its normalized
curvature derivatives in $\Delta'$ caps have supremum
$C_K\lambda$ and square-summed $L^2$ norm $C_K\sqrt D$, by
\eqref{eq:src16}, spatial Bernstein, and the bounded preceding-center
window.  These bounds hold for the fixed extra derivative orders
introduced below.

For one fixed cap-offset pattern, Fourier expansion of a smooth
rescaled multiplier produces integrals of spatially shifted products.
H\"older with a distinguished $L^2$ input, another $L^2$ input and
two $L^\infty$ inputs gives the base term in \eqref{eq:src21}.
For an earlier loss, use its distinguished $L^1$ input and three
$L^\infty$ inputs.  In summing the common cap vertex, Cauchy--Schwarz
handles the two $L^2$ sequences, or direct summation handles the
$L^1$ sequence.  There is no additional factor counting free caps.
Any further derivative falling on the distinguished input preserves
its gain, since $\Delta_m\leq\Delta'$.

\emph{The off-diagonal multiplier summation.}
On $|u|,|v|\lesssim_K\Delta'$ the divided symbol, with smooth
cutoffs, has bounded Fourier-kernel $L^1$ norm in its two rescaled
variables, and only boundedly many cap offsets occur.
For explicit rescaled bounds, write $s=u/\Delta'$, $t=v/\Delta'$,
$F_A(s)=1-e^{-cs^2/A^2}$, and
$E(s,t)=F_K(s)F_K(t)-F_1(s)F_1(t)$.
For the remaining regions, by symmetry let
$|v|\asymp Z\Delta'$ be the larger variable, with $Z$ dyadic and
large.  If $|u|\lesssim\Delta'/Z$, the once-divided symbol has kernel
cost $O_K(Z^{-2})$ at these derivative scales.  There are $O(Z)$
possible offsets for each cap vertex, giving total cost $O_K(Z^{-1})$.
In fact $E(s,t)=s^2e(s,t)$, with every required derivative of $e$
in coordinates $(Zs,t/Z)$ bounded uniformly in $Z$.
Consequently $E(s,t)/(2st)=(s/(2t))e(s,t)$ has fixed rescaled
derivative bounds $O_K(Z^{-2})$ after the cutoffs.  Anisotropic
rescaling preserves the $L^1$ norm of its Fourier kernel.

In the remaining dyadic regions $|u|\asymp U\Delta'$,
$1/Z\lesssim U\lesssim Z$, use \eqref{eq:frame-quartet} three
additional times.  After four divisions the symbol has kernel cost
\[
 O_K\bigl((U/Z)(UZ)^{-3}\bigr)
       =O_K(U^{-2}Z^{-4}).
\]
Indeed the undivided error and its fixed derivatives in coordinates
$(s/U,t/Z)$ cost $O_K(\min(U^2,1))\leq O_K(U^2)$, while each
division costs $(UZ)^{-1}$.
Smooth Gaussian derivatives give these estimates uniformly on
overlapping rectangles; the very small $u$ region uses a single
cutoff at scale $\Delta'/Z$.
The spatial sum constraint on the quartet leaves at most
$O(Z\max(1,U))$ offset patterns per fixed vertex.
Fixing the first cap leaves $O(\max(1,U))$ choices for the second
and $O(Z)$ for the fourth, and then $O(1)$ for the third by the sum
constraint.  The multiplier-induced spatial translations only
modulate Fourier factors and do not enlarge their supports, so they
do not change these counts.
Their combined cost is therefore
$O_K(Z^{-3}U^{-2}\max(1,U))$.
For fixed $Z$ the dyadic sum over $1/Z\lesssim U\leq1$ is
$O_K(Z^{-1})$, and that over $1\leq U\lesssim Z$ is
$O_K(Z^{-3})$.  Including the very small $u$ part and then summing
dyadic $Z$ converges.  At least the first curvature derivative remains
distinguished in every term of the expanded quartet identities.
The norm bounds above thus prove \eqref{eq:src21}, with no temporal
cap-support assumption and no dependence on the number of scales.

\emph{Telescoping the potential.}
At fixed scale $j$,
\begin{equation}\label{eq:frame-potential-loss}
 |\mathcal P_j(g_j)-\mathcal P_j(g_{j-1})|
 \lesssim_K\lambda^2l_j.
\end{equation}
Both Gaussian-filtered states have supremum $C_K\lambda$ by their
nearby compact cap pieces and rapid Gaussian overlap.  Also
\[
 \int_\R\|q_{j,a}R_j\|_1\frac{da}{\Delta_j}
 \lesssim\int_a\|d_a\|_1\lesssim_K\lambda^{-1}l_j.
\]
The first inequality follows by composing the Gaussian with the
compact synthesis projections: their $L^1$ kernels have rapidly
decaying overlap in the normalized difference of centers.
The inequality
$\bigl||u|^4-|v|^4\bigr|\lesssim(|u|+|v|)^3|u-v|$
then proves \eqref{eq:frame-potential-loss}.
Similarly $\mathcal P_n(g_n)\lesssim_K\lambda^2D$.
In summing \eqref{eq:frame-potential-increment}, insert and subtract
$\mathcal P_j(g_j)$ at each stage.  The remaining boundary terms are
$\mathcal P_n(g_n)-\mathcal P_0(g_0)$, with the initial term favorable.
The errors in \eqref{eq:src21} sum because
\[
 \sum_{j>m}(\Delta_m/\Delta_{j-1})^2\lesssim_K1,
 \qquad
 \sum_j(\Delta_0/\Delta_{j-1})^2\lesssim_K1,
 \qquad \sum_ml_m\leq D.
\]
This proves \eqref{eq:src22}.

All identities are global spacetime integrals.  There are finitely
many scales, and cap supremum bounds together with $L^2$ bounds give
finite quartic integrals.  Smooth localized approximations, if needed,
justify Fourier expansions; the established fixed derivative and
decay bounds pass to the limit.  No derivative in the center variable
$a$ is used.
\end{proof}

\begin{lemma}[Maximal recovery of the error fields]
\label{lem:frame-maximal}
The partial errors obey \eqref{eq:src19}.  If $\mathcal M$ is the
spacetime Hardy--Littlewood maximal operator, then also
\begin{equation}\label{eq:src23}
 |Y_k(z')|\lesssim
 (1+\Delta_k|z-z'|)^C\mathcal M(\sup_j|Y_j|)(z)+C_K\lambda^2.
\end{equation}
\end{lemma}

\begin{proof}
\emph{Recovering partial sums without a logarithm.}
Decompose the spatial outputs of $H_j$ and $W_j\Phi_j$ into smooth
Littlewood--Paley bands $r\asymp2^{-\ell}\Delta_j$, with $\ell$
bounded below by a fixed integer.  Lemma~\ref{lem:frame-flux} and
$\sum_jl_j\leq D$ bound the square sum of flux band norms by
$C_K\lambda^2D\,2^{-\ell}$.
For the growth term, evenness of $w$ at zero and its compact support
give, on the corresponding band,
\[
 |w(u/\Delta_j)-w(u/\Delta_{j-1})|^2
 \leq C_K2^{-2\ell}(G_j-G_{j-1})(u).
\]
Near zero the squared difference vanishes to fourth order, whereas
the Gaussian gap vanishes to second order; on the rest of the fixed
compact support their ratio is bounded.  Hence
\eqref{eq:src22} bounds the square sum of growth band norms by
$C_K\lambda^2D\,2^{-2\ell}$.

For a fixed $\ell$, successive stage bands have scale ratio $K$.
Choose $K$ larger than their fixed support ratio, so the spectra are
disjoint and ordered.  If $F_\ell$ is the sum over all stages in that
relative band, a smooth spatial low pass between successive bands
recovers each initial segment exactly and kills the later terms.
Its pointwise maximal value is bounded by a constant times the
spatial Hardy--Littlewood maximal function of $F_\ell$.
Plancherel for the disjoint spectra and the maximal $L^2$ bound,
integrated also in $t$, therefore bound maximal initial segments by
the square sum of their norms.  Summing $\ell$ by the triangle
inequality converges, with leakages $2^{-\ell/2}$ and $2^{-\ell}$.
Zero spatial frequency contributes no separate $L^2$ piece.
This proves \eqref{eq:src19}; no maximal inequality for arbitrary
orthogonal series is being used.

\emph{Local constancy.}
In each $H_j$ and $W_j\Phi_j$, spatial output support and the flux
expansion permit grouping into $O(1/\Delta_j)$ matched product bins
at bounded centers $b$.  The unprojected $d_a$ factors can again be
restricted to matched bins by the other input and output supports.
Every such product has
\[
 \| (\partial_t+2b\partial_x)^s(\text{product})\|_\infty
 \leq C_{s,K}\lambda^2\Delta_j^{2s}.
\]
This follows from \eqref{eq:src16}, the mask derivative bounds,
and bounded spatial $L^1$ kernels; the preceding caps per current
bin cost only a fixed factor depending on $K$.
Cut a product to affine temporal frequency
$|\tau+2b\xi|\lesssim\Delta_j$.  The smooth high pass divided by
the $s$th power of the rescaled affine frequency has integrable
kernel, for a fixed sufficiently large $s$.  The resulting error
is $C_{s,K}\lambda^2\Delta_j^s$ per bin, hence at most
$C_{s,K}\lambda^2\Delta_j^{s-1}$ for the stage.  These errors sum
geometrically for $s>1$.

The sum $\widetilde Y_k$ of the truncated products differs from
$Y_k$ by $O_K(\lambda^2)$ uniformly and has full spacetime Fourier
support in a ball of radius $O(\Delta_k)$: its spatial frequencies
are $O(\Delta_j)$ and its velocities are bounded.  A smooth spacetime
low pass at scale $\Delta_k$ recovers $\widetilde Y_k$, so it also
recovers $Y_k$ up to $O_K(\lambda^2)$.  Translating its Schwartz
kernel from $z'$ relative to a reference point $z$, and comparing
dyadic kernel annuli with the maximal function, proves
\eqref{eq:src23}.
\end{proof}

\Needspace{12\baselineskip}
\subsection{Stopping and sampling the two separated groups}

\begin{lemma}[A local bilinear sampling bound]
\label{lem:frame-local-bilinear}
Consider the two states from the preceding construction with
$0<\lambda_i\leq1$.
Let $Q$ be a square of side $\Delta_k^{-1}$ and let $\mathcal Z_Q\subset Q$
contain at most one designated point from each unit lattice square.
Suppose both states $g_{k-1}^{(i)}$ have modulus at least
$c\alpha_i$ at these points.  There is a Schwartz weight $\eta_Q$,
bounded below on $Q$ and with Fourier support of radius
$O(\Delta_k)$, for which
\begin{equation}\label{eq:frame-local-sampling}
 \alpha_1^2\alpha_2^2\#\mathcal Z_Q
 \lesssim_K\delta^{-C}\int_{\R^2}|\eta_Q(z')|^2
 \prod_{i=1}^2\left(W_{k-1}|g_{k-1}^{(i)}|^2(z')+1\right)\,dz',
\end{equation}
provided the ratio constants are sufficiently large.  The two
frequency intervals must retain a gap comparable to $\delta$.
\end{lemma}

\begin{proof}
Put $\Delta=\Delta_k$.  Partition each state by smooth spatial
lattice caps $Q_b$ of width $\Delta$.  Cut each piece to affine
modulation
\[
 |\tau+2b\xi-b^2|\lesssim\Delta.
\]
Its $D_b^r$ estimates show that this cutoff makes a uniform error
$C_{r,K}\lambda_i\Delta^r$ per cap, by the same inverse-derivative
high-pass argument as in Lemma~\ref{lem:frame-maximal}.
There are $O(1/\Delta)$ caps.  For a fixed sufficiently large $r$,
the total supremum error is $O_K(\lambda_i)$ and the sum of squared
errors is $O_K(\lambda_i^2)$.  Thus the truncated fields, denoted
$\widetilde g^{(i)}=\sum_bf_{i,b}$, still have modulus comparable
to $\alpha_i$ at the designated points.

Choose $\eta_Q$ as stated, with arbitrarily high fixed Schwartz
decay at the scale of $Q$.  Such a weight can be made by scaling a
function with compact smooth Fourier support and positive lower
bound on the unit square.  The full spacetime Fourier support of
$\eta_Q\widetilde g^{(1)}\widetilde g^{(2)}$ is bounded independently
of scales, since the frequency centers and velocities are bounded.
A Schwartz reproducing kernel at unit scale therefore bounds its
squared value at any point of a unit square by a rapidly weighted
local squared integral.  Summing over the distinct unit squares gives
\[
 \alpha_1^2\alpha_2^2\#\mathcal Z_Q
 \lesssim\|\eta_Q\widetilde g^{(1)}\widetilde g^{(2)}\|_2^2.
\]
This permits arbitrary designated points, including points on cell
closures, with only a fixed multiplicity adjustment.

The support of $f_{i,b}$ lies within $O(\Delta)$ of
$(b,-b^2)$ in spacetime frequency.  Hence the pair products after
multiplication by $\eta_Q$ have Fourier overlap bounded by a fixed
power of $\delta^{-1}$.  In detail, the map
\[
 (b,b')\longmapsto(b+b',b^2+(b')^2)
\]
has determinant $2(b'-b)$ on the two ordered separated intervals.
If two such mapped pairs are within $O(\Delta)$, their sums and
squared differences imply that their original ordered centers are
within $O(\delta^{-1}\Delta)$.  Since the cap lattice spacing is
$\Delta$, the overlap is at most $C\delta^{-C}$, uniformly in
the number of caps.  Plancherel then gives
\[
 \|\eta_Q\widetilde g^{(1)}\widetilde g^{(2)}\|_2^2
 \lesssim\delta^{-C}\int|\eta_Q|^2
                 \left(\sum_b|f_{1,b}|^2\right)
                 \left(\sum_{b'}|f_{2,b'}|^2\right).
\]

For completeness, the local spatial Bessel bound is
\[
 \sum_b|Q_bg(x,t)|^2
 \lesssim\int_\R(1+|y|)^{-P_0}|g(x-y/\Delta,t)|^2\,dy
 \lesssim_K W_{k-1}|g|^2(x,t).
\]
To obtain the first inequality, write the projection kernels as
modulations of a common Schwartz kernel.  The values $Q_bg(x,t)$
are Fourier coefficients on a fixed period of the periodization of
$y\mapsto\check\phi(y)g(x-y/\Delta,t)$.  Parseval and weighted
Cauchy--Schwarz for the periodization tails give the displayed
positive convolution.  The second inequality compares its scale
$\Delta^{-1}$ with $\Delta_{k-1}^{-1}=K\Delta^{-1}$.
The bounded squared sum of the temporal-cutoff errors consequently
gives
$\sum_b|f_{i,b}|^2\lesssim_KW_{k-1}|g_{k-1}^{(i)}|^2+1$,
proving \eqref{eq:frame-local-sampling}.
\end{proof}

\begin{proof}[Proof of Theorem~\ref{thm:frame}]
Apply Lemma~\ref{lem:frame-initial} and first work with one retained
layer of each field, normalized by its fine envelope bound.  Use the
same scales and masks above, with $\lambda_i=C_*/\alpha_i$ and
$D_i\lesssim H_i^*/B_i^2$.  The constants are chosen in this order:
the fixed profile and decay bounds, then $C_1$ and $K$ as in
Lemma~\ref{lem:frame-masks}, then a trigger level $C_0$, then $C_*$,
and finally the required lower bound for $\alpha_i$.
Choose $C_0$ sufficiently large relative to the uniform
$O_K(\lambda_i^2)$ errors in \eqref{eq:src23}, using
$\lambda_i\leq1$.

At a witness, until the first index $k$ at which either
$|Y_k^{(i)}|>C_0$, positivity in \eqref{eq:frame-recurrence} and
the witness baseline \eqref{eq:src14} give
\[
 \sum_{j<k}W_jL_j^{(i)}\lesssim1+C_0.
\]
By \eqref{eq:src18}, the total change in each field through those
stages, including the initial deletion, is at most
$C_K(1+C_0)\lambda_i^{-1}$.
Choose $C_*$ large enough that this is a small fixed fraction of
$\alpha_i$, since $\lambda_i^{-1}=\alpha_i/C_*$.
Both preceding states therefore retain modulus at least
$c\alpha_i$ at the witness.  A trigger must occur by stage $n$:
otherwise both final states retain those values, whereas the terminal
synthesis has only $O_K(1)$ normalized centers and hence supremum
$O_K(\lambda_i)$.  The last assertion uses the interval size
$O(\delta)$ and $\Delta_n\asymp_K c_0\delta$.

Use nested dyadic square grids of side $\Delta_k^{-1}$; this is
possible because $K$ and $\Delta_0$ are dyadic.  Each witness selects
the square of its first-trigger scale.  Take the maximal squares
$Q$ among these selected squares.  They are disjoint and cover all
witnesses; assign a witness to the unique maximal square containing
it.  If such a square has scale $k$, an assigned witness cannot have
triggered at a smaller index $j<k$: its larger dyadic square would
strictly contain $Q$ and contradict maximality.  Thus both
$g_{k-1}^{(i)}$ are large at every witness assigned to $Q$.
At least one witness in $Q$ triggers at index $k$.  Applying
\eqref{eq:src23} to that point yields, almost everywhere on $Q$,
\begin{equation}\label{eq:frame-trigger-maximal}
 \mathcal M_1(z)+\mathcal M_2(z)\gtrsim_K1,
 \qquad
 \mathcal M_i=\mathcal M\bigl(\sup_j|Y_j^{(i)}|\bigr).
\end{equation}

Apply Lemma~\ref{lem:frame-local-bilinear} to the assigned witnesses.
We explain how to bound its weighted integral using an arbitrary
reference point $z\in Q$.  Fix a witness $z_0\in Q$, common to both
fields.  The recurrence at stage $k-1$ gives
\[
 W_{k-1}|g_{k-1}^{(i)}|^2(z')
 \leq S_0^{0,(i)}(z')+|Y_{k-1}^{(i)}(z')|.
\]
The baseline \eqref{eq:src14}, and the fact that
$h\gtrsim\Delta_k^{-1}$ up to a fixed constant, bound its first
term by a fixed polynomial in
$1+\Delta_k\dist(z',Q)$.
For the second term use \eqref{eq:src23} at the reference point
$z$, noting $\Delta_{k-1}\leq\Delta_k$; if $k=1$, it is simply
zero.  We obtain
\[
 W_{k-1}|g_{k-1}^{(i)}|^2(z')+1
 \lesssim_K(1+\Delta_k\dist(z',Q))^{C}
                         (1+\mathcal M_i(z)).
\]
Choose the fixed Schwartz decay of $\eta_Q$ large enough to absorb
this polynomial.  Its integral scale is $|Q|$, and
$(1+\mathcal M_1)(1+\mathcal M_2)\lesssim
1+\mathcal M_1^2+\mathcal M_2^2$.  Thus the local sampling estimate
is at most
\[
 C_K\delta^{-C}|Q|
            (1+\mathcal M_1(z)^2+\mathcal M_2(z)^2)
\]
for almost every reference point $z\in Q$.  Average over $Q$ and
absorb the constant using \eqref{eq:frame-trigger-maximal}.  This proves
\[
 \alpha_1^2\alpha_2^2
 \#\{\text{witnesses assigned to }Q\}
 \lesssim_K\delta^{-C}\int_Q(\mathcal M_1^2+\mathcal M_2^2).
\]
Summing the disjoint squares, then applying the maximal $L^2$
inequality and \eqref{eq:src19}, gives
\[
 \#\lesssim_K\delta^{-C}\alpha_1^{-2}\alpha_2^{-2}
                  (\lambda_1^2D_1+\lambda_2^2D_2).
\]
Finally
$\lambda_i^2D_i\lesssim H_i^*/A_i^2$ in the original units.
The summable threshold splitting and excess amplitude decay in
Lemma~\ref{lem:frame-initial} allow summation over all initial layer
pairs.  This proves \eqref{eq:src8}.

The construction used only a fixed finite list of profile derivatives
and decay orders.  The derivative inductions are contractive, the
potential errors and relative bands sum geometrically, and all
coverings have fixed overlap.  Consequently no constant grows with
the number of stages or the array cardinality.  The only variable
separation loss is the fixed polynomial overlap bound in
Lemma~\ref{lem:frame-local-bilinear}.
\end{proof}

\begin{remark}[Use for composite tangential packets]
\label{rem:frame-composites}
The theorem uses only the stated derivative and phase-space bounds.
It therefore applies to a smoothly localized sum of smaller packets
treated as a single profile with amplitude equal to a bound for its
required derivatives.  Translating a time bin by $t_0$ replaces a
tangential intercept by $b+2\xi t_0$.  For a generic family its
phase-space counts remain bounded under this change when
$|t_0|\lesssim h^2$.  In the later composite-packet application,
larger shifts are also permitted because the carrier centers have
bounded count per $h^{-1}$ frequency bin and the intercept lattice
is translated separately at each fixed carrier.  It is this extra
spacing property that preserves the counts for those shifts.
On a normal slab of width $h$,
a Fourier expansion in the normalized normal variable produces mode
amplitudes $O_N(\langle n\rangle^{-N})$ times that derivative bound;
thresholds split with $\langle n\rangle^{-2}$ and sum as in the
initial layer reduction.  When omitted normal and time-cell
coordinates are at bounded relative distance, the tangential
envelope is bounded by a fixed multiple of the full cell envelope.
Thus this use introduces neither a dependence on the number of
component packets nor an additional count of slabs when their
coefficient energies have bounded overlap.
\end{remark}
